% =============================================================================
% Modern Professional Book Style Package
% Updated with Source Serif Pro and contemporary design
% =============================================================================

% Font encoding first
\usepackage[T1]{fontenc}
\usepackage[utf8]{inputenc}

% Microtype for superior typography
\usepackage[
activate={true,nocompatibility},
final,
tracking=true,
kerning=true,
spacing=true,
factor=1100,
stretch=10,
shrink=10
]{microtype}

% LANGUAGE SETTINGS ------------------------------------------------------
\usepackage[english,italian]{babel}
\addto\captionsitalian{%
	\renewcommand{\contentsname}{Indice}
	\renewcommand{\chaptername}{Capitolo}
}

% MODERN FONT SETUP - Source Serif Pro -----------------------------------
\usepackage[default,semibold]{sourceserifpro}
\usepackage{sourcesanspro}  % For headers and display text
\usepackage{sourcecodepro}  % For any code snippets

% Optimal line spacing for Source Serif
\linespread{1.18}

% TYPOGRAPHY ENHANCEMENTS ------------------------------------------------
\usepackage{lettrine}   % Drop caps
\usepackage{graphicx}   
\usepackage{xcolor}
\usepackage{tcolorbox} % Per creare riquadri personalizzabili
\tcbuselibrary{skins, breakable} % Librerie aggiuntive per i riquadri

% Modern color palette
\definecolor{modernblack}{RGB}{25, 25, 25}      % Softer than pure black
\definecolor{modernred}{RGB}{215, 38, 61}       % Contemporary accent
\definecolor{moderngray}{RGB}{110, 110, 110}    % For secondary elements
\definecolor{lightgray}{RGB}{245, 245, 245}     % For backgrounds
\definecolor{darkgray}{RGB}{70, 70, 70}         % For strong emphasis

% Professional paragraph settings
\setlength{\parskip}{0pt}                   
\setlength{\parindent}{1.2em}              
\frenchspacing
\widowpenalty=10000
\clubpenalty=10000
\raggedbottom  % Prevents ugly vertical stretching

% PAGE LAYOUT (17x24 cm - Saggio Grande) ---------------------------------
\usepackage[
paperwidth=170mm,
paperheight=240mm,
top=20mm,
bottom=25mm,
outer=20mm,
inner=25mm,
headsep=8mm,
footskip=12mm
]{geometry}

% Better line breaking
\emergencystretch=3em
\hfuzz=.5pt  % Suppress trivial overfull warnings

% TABLE OF CONTENTS STYLING ----------------------------------------------
\usepackage[titles]{tocloft}
\renewcommand{\cftchapfont}{\sffamily\normalsize}
\renewcommand{\cftchappagefont}{\sffamily\normalsize\color{moderngray}}
\renewcommand{\cftchapleader}{\cftdotfill{\cftdotsep}}
\renewcommand{\cftdotsep}{2.0}
\setlength{\cftbeforechapskip}{0.5em}

% CHAPTER STYLING --------------------------------------------------------
\usepackage{titlesec}

% Modern chapter style with sans-serif
\titleformat{\chapter}[display]
{\normalfont\sffamily\centering}
{\huge\textls[50]{\MakeUppercase{\chaptertitlename}}\quad\thechapter}
{25pt}
{\Huge\bfseries}

\titlespacing*{\chapter}{0pt}{40pt}{60pt}

% Section and subsection styles
\titleformat{\section}
{\normalfont\sffamily\Large\bfseries\color{darkgray}}
{\thesection}
{1em}
{}

\titleformat{\subsection}
{\normalfont\sffamily\large\bfseries\color{darkgray}}
{\thesubsection}
{1em}
{}

\titleformat{\subsubsection}
{\normalfont\sffamily\normalsize\bfseries\color{moderngray}}
{\thesubsubsection}
{1em}
{}

% HEADERS AND FOOTERS ---------------------------------------------------
\usepackage{fancyhdr}
\setlength{\headheight}{25pt}

% Main matter page style
\fancypagestyle{mainmatter}{%
	\fancyhf{}
	% Even pages: page number and author
	\fancyhead[LE]{\sffamily\footnotesize\color{moderngray}\thepage\quad\textls[50]{\MakeUppercase{\authorname}}}
	% Odd pages: chapter title and page number
	\fancyhead[RO]{\sffamily\footnotesize\color{moderngray}\textls[50]{\MakeUppercase{\leftmark}}\quad\thepage}
	% No header rule for cleaner look
	\renewcommand{\headrulewidth}{0pt}
}

% Plain style for chapter openings
\fancypagestyle{plain}{%
	\fancyhf{}
	\fancyfoot[C]{\sffamily\small\color{moderngray}\thepage}
	\renewcommand{\headrulewidth}{0pt}
}

% QUOTATION ENVIRONMENT --------------------------------------------------
\usepackage{csquotes}
\renewenvironment{quotation}
{\list{}{\listparindent 0em%
		\itemindent    0em
		\leftmargin    2em
		\rightmargin   2em
		\parsep        0.5em
		\topsep        0.5em}%
	\item\relax\small\color{darkgray}}
{\endlist}

% Block quotes with modern style
\newenvironment{modernquote}
{\begin{quotation}\sffamily\itshape}
	{\end{quotation}}

% LISTS - Modern Spacing -------------------------------------------------
\usepackage{enumitem}
\setlist[itemize]{
	topsep=0.5em,
	itemsep=0.25em,
	parsep=0.25em,
	leftmargin=1.5em,
	label={\color{modernred}\textbullet}
}
\setlist[enumerate]{
	topsep=0.5em,
	itemsep=0.25em,
	parsep=0.25em,
	leftmargin=1.5em
}

% CAPTIONS ---------------------------------------------------------------
\usepackage[
font={scriptsize,sf},
labelfont={bf,color=darkgray},
textfont={color=moderngray},
justification=centering,
margin=0.5cm
]{caption}
\setlength{\abovecaptionskip}{10pt}
\setlength{\belowcaptionskip}{10pt}

% FOOTNOTES --------------------------------------------------------------
\usepackage[bottom,hang,flushmargin]{footmisc}
\setlength{\footnotesep}{12pt}
\renewcommand{\footnotelayout}{\footnotesize\color{moderngray}}

% CUSTOM COMMANDS --------------------------------------------------------
% First few words in sans-serif bold
\newcommand{\firstwords}[1]{{\sffamily\bfseries #1}}

% Small caps with tracking
\newcommand{\smallcaps}[1]{\textsc{\textls[50]{#1}}}

% Book title formatted
\newcommand{\booktitleformatted}{{\sffamily\textit{\booktitle}}}

% Chapter quotes
\newcommand{\chapterquote}[2]{%
	\vspace{-0.5em}
	\begin{flushright}
		\begin{minipage}{0.8\textwidth}
			\raggedleft
			\sffamily\small\itshape\color{moderngray}
			``#1''\\[0.5em]
			\sffamily\footnotesize\color{darkgray}--- #2
		\end{minipage}
	\end{flushright}
	\vspace{1em}
}

% Modern lettrine settings
\renewcommand{\LettrineFontHook}{\sffamily\bfseries\color{modernblack}}
\setlength{\DefaultNindent}{0pt}
\setlength{\DefaultSlope}{0pt}
\renewcommand{\LettrineTextFont}{\sffamily}
\setlength{\DefaultFindent}{0.15em}

% PAGE NUMBERS WITH STYLE ------------------------------------------------
\usepackage{lastpage}

% HYPERREF (load last) ---------------------------------------------------
\usepackage[
colorlinks=true,
linkcolor=modernblack,
citecolor=modernred,
urlcolor=modernred,
bookmarks=true,
bookmarksopen=true,
bookmarksnumbered=true,
pdfauthor={\authorname},
pdftitle={\booktitle},
pdfsubject={Neuroscienze e Tecnologia},
pdfkeywords={cervello, tecnologia, evoluzione, psicologia},
pdfproducer={LaTeX with Source Serif Pro},
pdfcreator={pdfLaTeX}
]{hyperref}

% URL STYLE
\urlstyle{same}  % URLs in same font as text

% ADDITIONAL MODERN TOUCHES ----------------------------------------------
% Better spacing around display math (if used)
\usepackage{amsmath}
\AtBeginDocument{%
	\setlength\abovedisplayskip{12pt plus 3pt minus 3pt}%
	\setlength\belowdisplayskip{12pt plus 3pt minus 3pt}%
}

% Epigraphs for chapter quotes
\usepackage{epigraph}
\setlength{\epigraphwidth}{0.7\textwidth}
\renewcommand{\epigraphflush}{flushleft}
\renewcommand{\sourceflush}{flushleft}

% Optional: Add TikZ for any diagrams
\usepackage{tikz}
\usetikzlibrary{shapes,arrows,positioning}

% Optional: Better tables
\usepackage{booktabs}
\renewcommand{\arraystretch}{1.2}
% =============================================================================
% INDEXING PACKAGE AND CONFIGURATION
% =============================================================================
\usepackage[xindy]{imakeidx}

% Stile professionale per l'indice
\indexsetup{
	othercode=\small,
	firstpagestyle=plain,
	headers={Indice analitico}{Indice analitico}
}

% Indice analitico principale
\makeindex[
name=main,
title=Indice analitico,
program=texindy,
options=-L italian -C utf8 -I latex,
columns=2 
]

% Indice dei bias cognitivi
\makeindex[
name=cognitive,
title=Indice dei bias cognitivi,
program=texindy,
options=-L italian -C utf8 -I latex,
columns=2 
]

% =============================================================================
% STILI PER RIQUADRI CONCETTUALI E CODICE
% =============================================================================

% --- Riquadro per concetti chiave ---
\newtcolorbox{concetto}[1][]{
	colback=lightgray,       
	colframe=darkgray,       
	fonttitle=\sffamily\bfseries,
	coltitle=modernblack,
	breakable,
	enhanced,
	arc=3mm,
	boxrule=0.8pt,
	title=#1
}

% --- Riquadro per lo pseudo-codice ---
\newtcolorbox{codice}[1][]{
	colback=darkgray,       
	coltext=white,
	fontupper=\ttfamily,      
	breakable,
	arc=1mm,
	boxrule=0pt,
	title=#1,
	fonttitle=\sffamily\bfseries,
	coltitle=white
}

% --- Riquadro per Dati e Grafici ---
\newtcolorbox{boxdati}[1][]{
	colback=white,            
	colframe=modernred,       
	fonttitle=\sffamily\bfseries,
	coltitle=white,
	breakable,
	enhanced,
	arc=3mm,
	boxrule=0.8pt,
	title=#1
}

% =============================================================================
% STILI PER DIALOGO UMANO-AI
% =============================================================================

\definecolor{humancolor}{RGB}{25, 25, 25}      
\definecolor{aicolor}{RGB}{70, 130, 180}       

\newcommand{\humanbox}[1]{%
	\vspace{0.8em}%
	\begin{tcolorbox}[
		colback=white,
		colframe=humancolor,
		leftrule=3pt,
		rightrule=0pt,
		toprule=0pt,
		bottomrule=0pt,
		left=10pt,
		right=5pt,
		top=5pt,
		bottom=5pt
		]
		\noindent\textbf{\sffamily\color{humancolor}UMANO:} #1
	\end{tcolorbox}%
	\vspace{0.3em}%
}

\newcommand{\aibox}[1]{%
	\vspace{0.8em}%
	\begin{tcolorbox}[
		colback=aicolor!5,
		colframe=aicolor,
		leftrule=3pt,
		rightrule=0pt,
		toprule=0pt,
		bottomrule=0pt,
		left=10pt,
		right=5pt,
		top=5pt,
		bottom=5pt
		]
		\noindent\textbf{\sffamily\color{aicolor}AI:} #1
	\end{tcolorbox}%
	\vspace{0.3em}%
}